\documentclass[10pt,a4paper]{amsart}
\usepackage[margin=19mm]{geometry}
\usepackage{amsmath,amssymb,mathtools}
\usepackage[scr=rsfs,cal=euler]{mathalfa}
\usepackage{xfrac}
\usepackage[unicode,hidelinks,bookmarks=false]{hyperref}
\newcommand{\ie}{i.e.\ }
\newcommand{\cf}{cf.\ }
\newcommand{\resp}{respectively\ }
\newcommand{\Subsection}{\subsection}
\providecommand{\nobreakdash}{\nobreak\hbox{-}}
\setlength{\emergencystretch}{2em}
\multlinegap0pt
\usepackage{xcolor}
\newtheorem{lem}{Lemma}
\title{Sharp Fourier inequalities and lattice point discrepancy for $l^p$-balls}
\author{Martin Lind}
\date{Source: arXiv:2208.07837v2 (CC BY 4.0)}
\begin{document}
\maketitle

\noindent\emph{Source and licence.} Sharp Fourier inequalities and lattice point discrepancy for $l^p$-balls. Version arXiv:2208.07837v2 (CC BY 4.0), distributed by arXiv under the Creative Commons Attribution 4.0 International licence, http://creativecommons.org/licenses/by/4.0/, as recorded in the arXiv licence metadata for this exact version and shown on the abstract page https://arxiv.org/abs/2208.07837v2. Full licence notice: \texttt{licenses/arxiv-2208.07837-CC-BY-4.0.txt}.\par\smallskip

\noindent\textit{Editorial scope.} Counted unit: the article's lem x-p-lem (source0.tex lines 329--346). The article's complete original proof of it is delivered with it (source0.tex lines 347--398, one \texttt{proof} environment of 52 lines). No mathematics is written, repaired or re-derived: the statement and the proof are the article's own lines, in the article's own notation, and no retained line is substituted or re-flowed.  Retained uncounted as the source-local context the proof needs: lines 130--133.  Nothing was omitted from any retained span, so the package carries no omission record.  No retained line contains a citation, so the wrapper carries no bibliography and no bibliography entry.  No retained line contains a figure environment, an image-inclusion command or drawing code, so no image is delivered and the package declares no asset dependency.  Editorial formatting only, in addition: an AMS \texttt{amsart} preamble that restates the article's own declarations and macros used by the retained lines, namely the environments \texttt{lem} on separate counters, together with the standard packages those lines need.  The retained environments print in this document as Lemma 1 in order of appearance, whereas the article prints the same items with the numbering of its own full document.  The complete native source of the article as distributed in the arXiv e-print for this exact version is bundled unchanged as \texttt{sources/129537/source0.tex}.
\begin{equation}
\label{phi-Function}
\varphi_p(x)=\left(1-x^p\right)^{1/p}.
\end{equation}

\begin{lem}
\label{x-p-lem}
Let $\varphi_p$ be defined by (\ref{phi-Function}). The function $|\varphi_p''|$ has a unique minimum $x^*= x^*(p)$ on $[0,1]$. The minimum value satisfies
\begin{equation}
    \label{phi-p2ndDerivMin1}
    |\varphi_p''(x^*)|=\min_{0\le x\le1}|\varphi_p''(x)|=(p-1)m(p),
\end{equation}
where $m(p)$ is a decreasing function on $[1,2]$ with $m(1)=4$, $m(2)=1$. Further,
\begin{equation}
    \label{x-p-lem1}
    \frac{-1}{\varphi_p'(x^*)}\ge1,
\end{equation}
and 
\begin{equation}
    \label{x-p-lem2}
    \lim_{p\rightarrow1+}\frac{-1}{\varphi_p'(x^*)}=1.
\end{equation}
\end{lem}

\begin{proof}
We have
\begin{equation}
    \label{1stDervi-phi-p}
    \varphi_p'(x)=-x^{p-1}(1-x^p)^{1/p-1},
\end{equation}
\begin{equation}
    \label{2ndDervi-phi-p}
    \varphi_p''(x)=-(p-1)x^{p-2}(1-x^p)^{1/p-2},
\end{equation}
and
\begin{equation}
    \label{3rdDervi-phi-p}
    \varphi_p^{(3)}(x)=-(p-1)x^{p-3}(1-x^p)^{1/p-3}\left(x^p(p+1)+p-2\right).
\end{equation}
It is clear that the third derivative (\ref{3rdDervi-phi-p}) has only the zero
\begin{equation}
    \label{x-star}
    x^*\equiv x^*(p)=\left(\frac{2-p}{p+1}\right)^{1/p}
\end{equation}
on $(0,1)$. Further, $\varphi_p^{(3)}(x)>0$ for $x<x^*$ and $\varphi_p^{(3)}(x)<0$ for $x>x^*$. Hence, $\varphi_p''$ has a unique interior maximum at $x^*$. Since $\varphi_p''$ is strictly negative, it follows that $|\varphi_p''|$ has a unique interior minimum at $x^*$.
Inserting the expression (\ref{x-star}) into (\ref{2ndDervi-phi-p}) and taking absolute value yields
$$
|\varphi_p''(x^*)|=(p-1)(2-p)^{1-2/p}(2p-1)^{1/p-2}(p+1)^{1+1/p}.
$$
Define
$$
m(p)=(2-p)^{1-2/p}(2p-1)^{1/p-2}(p+1)^{1+1/p}\quad p\in[1,2), 
$$
and
$$
m(2)=\lim_{p\rightarrow2-}m(p).
$$
Clearly $m(1)=4$. By logarithmic differentiation,
$$
\frac{dm}{dp}=\frac{m(p)}{p^2}(-\ln(p+1)-\ln(2p-1)+2\ln(2-p))<0
$$
for $p\in(1,2)$. To demonstrate $m(2)=1$, we calculate the limit defining $m(2)$:
\begin{eqnarray}
\nonumber
\lim_{p\rightarrow2-}m(p)&=&\lim_{p\rightarrow2-}(2-p)^{1-2/p}\times 3^{-3/2}\times3^{3/2}\\
\nonumber
&=&\lim_{p\rightarrow2-}\exp\left(\left(1-\frac{2}{p}\right)\ln(2-p)\right)\\
\nonumber
&=&\lim_{t\rightarrow0+}\exp\left(\frac{-t\ln(t)}{(2-t)}\right)=\exp(0)=1.
\end{eqnarray}
To prove (\ref{x-p-lem1}), we insert the expression (\ref{x-star}) into (\ref{1stDervi-phi-p}) and get
$$
\varphi_p'(x^*)=-\left(\frac{2-p}{2p-1}\right)^{1-1/p}\quad\Leftrightarrow\quad\frac{-1}{\varphi_p'(x^*)}=\left(\frac{2p-1}{2-p}\right)^{1-1/p}.
$$
Since $p\ge1$, we have $(2p-1)/(2-p)\ge1$, thus proving (\ref{x-p-lem1}). Finally, (\ref{x-p-lem2}) follows by direct evaluation.
\end{proof}

\end{document}
